% Part: lambda-calculus
% Chapter: syntax
% Section: beta

\documentclass[../../../include/open-logic-section]{subfiles}

\begin{document}

\olfileid{lam}{syn}{bet}
\olsection{اختزال~$\beta$}

إذا تأملنا~$(\lambd[\OLId{latin-ordinary}{m}][(\lambd[\OLId{latin-ordinary}{y}][\OLId{latin-ordinary}{y}]) \OLId{latin-ordinary}{m}])$ ظهر لنا أن بينه وبين~$\lambd[\OLId{latin-ordinary}{m}][\OLId{latin-ordinary}{m}]$ صلة:
فالثاني ينبغي أن ينتج من تبسيط الأول.
ومفهوم \emph{اختزال~$\beta$} هو الصياغة الصورية لهذا الحدس.

\begin{defn}[تقلّص~$\beta$، $\bredone$] \ollabel{defn:betacontr}
  علاقة \emph{تقلّص~$\beta$} ($\bredone$) أصغر علاقة متوافقة على الحدود تحقق
  \[
    (\lambd[\OLId{latin-ordinary}{x}][\OLId{latin-looped}{N}])\OLId{latin-looped}{Q} \bredone \Subst{\OLId{latin-looped}{N}}{\OLId{latin-looped}{Q}}{\OLId{latin-ordinary}{x}} 
  \]
  ونقول إن~$\OLId{latin-looped}{P}$ \emph{يتقلّص بحسب~$\beta$} إلى~$\OLId{latin-looped}{Q}$ متى كان~$\OLId{latin-looped}{P} \bredone
  \OLId{latin-looped}{Q}$.
  وكل حد على صورة~$(\lambd[\OLId{latin-ordinary}{x}][\OLId{latin-looped}{N}])\OLId{latin-looped}{Q}$ يسمى \emph{موضع اختزال}.
\end{defn}

\begin{prob} \ollabel{prob:def}
  بيّن تفصيلًا التعريفات الاستقرائية المكافئة لتقلّص~$\beta$،
  على مثال ما صنعناه لتغيير المتغير المقيد في~\olref[lam][syn][alp]{defn:aconvone}.
\end{prob}
  
\begin{defn}[اختزال~$\beta$، $\bred$] \ollabel{defn:betared}
  علاقة \emph{اختزال~$\beta$} ($\bred$) أصغر علاقة انعكاسية متعدية على الحدود
  تشتمل على~$\bredone$. ونقول إن~$\OLId{latin-looped}{P}$ يختزل بحسب~$\beta$ إلى
  $\OLId{latin-looped}{Q}$ متى كان~$\OLId{latin-looped}{P} \bred \OLId{latin-looped}{Q}$.
\end{defn}

وحيث لا يلتبس المراد، نستعمل~$\redone$ مكان~$\bredone$،
و$\red$ مكان~$\bred$.

فالمعنى غير الصوري لـ$\OLId{latin-looped}{M} \bred \OLId{latin-looped}{N}$ هو إمكان الانتقال من~$\OLId{latin-looped}{M}$ إلى
$\OLId{latin-looped}{N}$ بصفر من خطوات تقلّص~$\beta$ أو بعدة خطوات منه.

\begin{defn}[السواء بحسب~$\beta$]
الحد الذي لا يقبل مزيدًا من التقلّص بحسب~$\beta$ يسمى \emph{سويًا بحسب~$\beta$}.
\end{defn}

إذا تحقق~$\OLId{latin-looped}{M} \bred \OLId{latin-looped}{N}$ وكان~$\OLId{latin-looped}{N}$ سويًا بحسب~$\beta$، سمينا~$\OLId{latin-looped}{N}$
\emph{صورة سوية} للحد~$\OLId{latin-looped}{M}$. وهل يمكن أن تتعدد الصورة السوية لحد واحد؟
سنرى لاحقًا أنها، إن وجدت، وحيدة.

ونبيّن هذه المعاني بالأمثلة:
\begin{enumerate}
\item نحصل على الاختزال
\begin{align*}
(\lambd[\OLId{latin-ordinary}{x}][xxy]) \lambd[\OLId{latin-ordinary}{z}][\OLId{latin-ordinary}{z}] & \redone (\lambd[\OLId{latin-ordinary}{z}][\OLId{latin-ordinary}{z}])(\lambd[\OLId{latin-ordinary}{z}][\OLId{latin-ordinary}{z}]) \OLId{latin-ordinary}{y} \\
& \redone (\lambd[\OLId{latin-ordinary}{z}][\OLId{latin-ordinary}{z}]) \OLId{latin-ordinary}{y} \\
& \redone \OLId{latin-ordinary}{y}
\end{align*}
\item وقد يزيد ما نسميه تبسيطًا في تعقيد الحد:
\begin{align*}
(\lambd[\OLId{latin-ordinary}{x}][xxy])(\lambd[\OLId{latin-ordinary}{x}][xxy]) & \redone (\lambd[\OLId{latin-ordinary}{x}][xxy])(\lambd[\OLId{latin-ordinary}{x}][xxy])\OLId{latin-ordinary}{y} \\
& \redone (\lambd[\OLId{latin-ordinary}{x}][xxy])(\lambd[\OLId{latin-ordinary}{x}][xxy])yy \\
& \redone \dots
\end{align*}
\item وقد يبقى الحد بعده على حاله:
\[
(\lambd[\OLId{latin-ordinary}{x}][xx])(\lambd[\OLId{latin-ordinary}{x}][xx]) \redone (\lambd[\OLId{latin-ordinary}{x}][xx])(\lambd[\OLId{latin-ordinary}{x}][xx])
\]
\item وقد تتعدد طرق الاختزال؛ فمن ذلك
\[
(\lambd[\OLId{latin-ordinary}{x}][(\lambd[\OLId{latin-ordinary}{y}][yx]) \OLId{latin-ordinary}{z}]) \OLId{latin-ordinary}{v} \redone (\lambd[\OLId{latin-ordinary}{y}][yv]) \OLId{latin-ordinary}{z}
\]
إذا قلصنا التطبيق الخارجي أولًا، وكذلك
\[
(\lambd[\OLId{latin-ordinary}{x}][(\lambd[\OLId{latin-ordinary}{y}][yx]) \OLId{latin-ordinary}{z}]) \OLId{latin-ordinary}{v} \redone (\lambd[\OLId{latin-ordinary}{x}][zx]) \OLId{latin-ordinary}{v}
\]
إذا قلصنا التطبيق الداخلي أولًا. ومع هذا الاختلاف، يختزل الناتجان في هذا المثال إلى حد واحد،
هو~$zv$.
\end{enumerate}

وليست وحدة المآل في المثال الأخير اتفاقًا عارضًا؛ بل هي شاهد على خاصية عميقة مهمة لحساب لامبدا،
تسمى خاصية تشيرش--روسر.

\begin{digress}
  لما كان اختزال الحد بحسب~$\beta$ يجري، في العموم، بطرق مختلفة،
  وضعت له استراتيجيات كثيرة؛ وأشهرها \emph{الاستراتيجية الطبيعية}.
  وهي تقلّص دائمًا موضع الاختزال \emph{الأبعد إلى اليسار}،
  وتعين موضعه بنقطة ابتداء كتابته في الحد.
  ومن منفعتها أن إمكان الوصول إلى صورة سوية باستراتيجية ما يكافئ إمكان الوصول إليها بالاستراتيجية الطبيعية.
  وهذه هي الاستراتيجية التي نتبعها فيما يأتي، إلا إذا نصصنا على خلافها.
\end{digress}

\begin{defn}[التكافؤ بحسب~$\beta$، $\equal$]
  علاقة \emph{التكافؤ بحسب~$\beta$} ($\equal$) تعرّف استقرائيًا بالقواعد الآتية:
  \begin{enumerate}
  \item $\OLId{latin-looped}{M} \equal \OLId{latin-looped}{M}$.
  \item من~$\OLId{latin-looped}{M} \equal \OLId{latin-looped}{N}$ يلزم~$\OLId{latin-looped}{N} \equal \OLId{latin-looped}{M}$.
  \item من~$\OLId{latin-looped}{M} \equal \OLId{latin-looped}{N}$ و$\OLId{latin-looped}{N} \equal \OLId{latin-looped}{O}$ يلزم~$\OLId{latin-looped}{M} \equal \OLId{latin-looped}{O}$.
  \item من~$\OLId{latin-looped}{M} \equal \OLId{latin-looped}{N}$ يلزم~$PM \equal PN$.
  \item من~$\OLId{latin-looped}{M} \equal \OLId{latin-looped}{N}$ يلزم~$MQ \equal NQ$.
  \item من~$\OLId{latin-looped}{M} \equal \OLId{latin-looped}{N}$ يلزم~$\lambd[\OLId{latin-ordinary}{x}][\OLId{latin-looped}{M}] \equal \lambd[\OLId{latin-ordinary}{x}][\OLId{latin-looped}{N}]$.
  \item $(\lambd[\OLId{latin-ordinary}{x}][\OLId{latin-looped}{N}])\OLId{latin-looped}{Q} \equal \Subst{\OLId{latin-looped}{N}}{\OLId{latin-looped}{Q}}{\OLId{latin-ordinary}{x}}$.
  \end{enumerate}
\end{defn}

فبالقواعد الثلاث الأولى تكون العلاقة تكافؤًا، وبالثلاث التالية تكون متوافقة،
وبالقاعدة الأخيرة تشتمل على تقلّص~$\beta$.

والمعنى غير الصوري أن تكافؤ حدين بحسب~$\beta$ يكافئ إمكان الانتقال من أحدهما إلى الآخر
بصفر أو أكثر من خطوات تقلّص~$\beta$ أو خطوات معكوس تقلّص~$\beta$.
ونعني بمعكوس تقلّص~$\beta$ أن~$\OLId{latin-looped}{M}$ يتقلّص عكسيًا بحسب~$\beta$ إلى~$\OLId{latin-looped}{N}$
إذا وفقط إذا تقلّص~$\OLId{latin-looped}{N}$ بحسب~$\beta$ إلى~$\OLId{latin-looped}{M}$.

وسنضيف إلى هذه القواعد قواعد أخرى توسع العلاقة.
ونرمز إلى التكافؤ الناتج بـ$\equal[\OLId{latin-looped}{X}]$، حيث~$\OLId{latin-looped}{X}$ القاعدة المضافة.

\end{document}

